\documentclass[10pt,a4paper]{amsart}
\usepackage{amsmath,amssymb,amsthm,graphicx}
\usepackage{verbatim}
\newenvironment{reference}{\comment}{\endcomment}
%\newenvironment{reference}{}{}
\newenvironment{slogan}{\comment}{\endcomment}
\newenvironment{history}{\comment}{\endcomment}

% For commutative diagrams we use Xy-pic
\usepackage[all]{xy}

% We use 2cell for 2-commutative diagrams.
\xyoption{2cell}
% Theorem environments.
%
\theoremstyle{plain}
\newtheorem{theorem}[subsection]{Theorem}
\newtheorem{proposition}[subsection]{Proposition}
\newtheorem{lemma}[subsection]{Lemma}

\theoremstyle{definition}
\newtheorem{definition}[subsection]{Definition}
\newtheorem{example}[subsection]{Example}
\newtheorem{exercise}[subsection]{Exercise}
\newtheorem{situation}[subsection]{Situation}

\theoremstyle{remark}
\newtheorem{remark}[subsection]{Remark}
\newtheorem{remarks}[subsection]{Remarks}

\numberwithin{equation}{subsection}

% Macros
%
\def\lim{\mathop{\mathrm{lim}}\nolimits}
\def\colim{\mathop{\mathrm{colim}}\nolimits}
\def\Spec{\mathop{\mathrm{Spec}}}
\def\Hom{\mathop{\mathrm{Hom}}\nolimits}
\def\Ext{\mathop{\mathrm{Ext}}\nolimits}
\def\SheafHom{\mathop{\mathcal{H}\!\mathit{om}}\nolimits}
\def\SheafExt{\mathop{\mathcal{E}\!\mathit{xt}}\nolimits}
\def\Sch{\mathit{Sch}}
\def\Mor{\mathop{\mathrm{Mor}}\nolimits}
\def\Ob{\mathop{\mathrm{Ob}}\nolimits}
\def\Sh{\mathop{\mathit{Sh}}\nolimits}
\def\NL{\mathop{N\!L}\nolimits}
\def\CH{\mathop{\mathrm{CH}}\nolimits}
\def\proetale{{pro\text{-}\acute{e}tale}}
\def\etale{{\acute{e}tale}}
\def\QCoh{\mathit{QCoh}}
\def\Ker{\mathop{\mathrm{Ker}}}
\def\Im{\mathop{\mathrm{Im}}}
\def\Coker{\mathop{\mathrm{Coker}}}
\def\Coim{\mathop{\mathrm{Coim}}}

% Boxtimes
%
\DeclareMathSymbol{\boxtimes}{\mathbin}{AMSa}{"02}

%
% Macros for moduli stacks/spaces
%
\def\QCohstack{\mathcal{QC}\!\mathit{oh}}
\def\Cohstack{\mathcal{C}\!\mathit{oh}}
\def\Spacesstack{\mathcal{S}\!\mathit{paces}}
\def\Quotfunctor{\mathrm{Quot}}
\def\Hilbfunctor{\mathrm{Hilb}}
\def\Curvesstack{\mathcal{C}\!\mathit{urves}}
\def\Polarizedstack{\mathcal{P}\!\mathit{olarized}}
\def\Complexesstack{\mathcal{C}\!\mathit{omplexes}}
% \Pic is the operator that assigns to X its picard group, usage \Pic(X)
% \Picardstack_{X/B} denotes the Picard stack of X over B
% \Picardfunctor_{X/B} denotes the Picard functor of X over B
\def\Pic{\mathop{\mathrm{Pic}}\nolimits}
\def\Picardstack{\mathcal{P}\!\mathit{ic}}
\def\Picardfunctor{\mathrm{Pic}}
\def\Deformationcategory{\mathcal{D}\!\mathit{ef}}
\usepackage{cleveref}
\crefname{originallemma}{original lemma}{original lemmas}
\Crefname{originallemma}{Original lemma}{Original lemmas}
\makeatletter
\newlabel{perfect-lemma-quasi-coherence-internal-hom}{{0A6H}{source}{}{sourcepointer-1}{}}
\newlabel{perfect-lemma-quasi-coherence-internal-hom@cref}{{[originallemma][0][]0A6H}{[1][0][]source}{}{}{}}
\makeatother
\crefname{originallemma}{original lemma}{original lemmas}
\Crefname{originallemma}{Original lemma}{Original lemmas}
\makeatletter
\newlabel{dualizing-lemma-localization-injective-modules}{{0A6I}{source}{}{sourcepointer-2}{}}
\newlabel{dualizing-lemma-localization-injective-modules@cref}{{[originallemma][0][]0A6I}{[1][0][]source}{}{}{}}
\makeatother
\crefname{originallemma}{original lemma}{original lemmas}
\Crefname{originallemma}{Original lemma}{Original lemmas}
\makeatletter
\newlabel{algebra-lemma-hom-from-finitely-presented}{{0583}{source}{}{sourcepointer-3}{}}
\newlabel{algebra-lemma-hom-from-finitely-presented@cref}{{[originallemma][0][]0583}{[1][0][]source}{}{}{}}
\makeatother
\title{Internal Hom on an affine scheme computes the algebraic Hom for a coherent first argument}
\author{The Stacks Project Authors}
\date{stacks-project:duality.tex:local-snapshot-20261012; GNU FDL 1.2 or later}
\begin{document}
\maketitle
\noindent\textit{Editorial note.} All delivered mathematical text is the author's own, copied from the pinned Stacks Project chapter file duality.tex; no statement or proof step was rewritten, reordered or removed.\par\medskip
\noindent\textit{Editorial note.} Preamble provenance: the pinned chapter file opens with an input of the shared preamble and carries no preamble text of its own, so the project's single Stacks-wide preamble staged at sources/stacks-duality/preamble.tex is excerpted verbatim under the recipe fixed by LANE\_BRIEF 22.1, as three spans, lines 8-13 (verbatim package and the reference/slogan/history environments), lines 14-18 (Xy-pic setup) and lines 154-217 (the native theorem declarations and mathematical macros). Nothing was added, removed or redefined; the wrapper class and the amsmath/amssymb/amsthm/graphicx base packages are supplied by the assembler.\par\medskip
\noindent\textit{Editorial note.} The counted statement and proof contain no cross reference, no citation and no included file, so the delivered body is self contained.\par\medskip
\noindent\textit{Editorial note.} Canonical location of the counted proposition: Stacks tag 0A88 (sources/stacks-duality/pinned-original-source.tex line 366).\par\medskip
\noindent\textit{Allowed standard prerequisites (editorial).}\begin{itemize}
\item Cohomology sheaves of the internal Hom on an affine scheme. Derived Categories of Schemes, Lemma labelled perfect-lemma-quasi-coherence-internal-hom, cited at proof lines 386-390. Identifies the presheaf whose sheafification is computed.
\item Localisation of injective modules over a Noetherian ring. Dualizing Complexes, Lemma labelled dualizing-lemma-localization-injective-modules, cited at proof lines 396-399. Makes I otimes\_A A\_f injective.
\item Hom from a finitely presented module commutes with localisation. Algebra, Lemma labelled algebra-lemma-hom-from-finitely-presented, cited at proof lines 406-410. Gives the last equality of the computation.
\end{itemize}\medskip
\ifdefined\hypertarget\hypertarget{sourcepointer-1}{}\fi
\noindent\textit{Source reference.} Original lemma 0A6H: sources/stacks-chow/tags.txt, source lines 12845--12845. Cross-chapter standard prerequisite located in the project tag list staged in the local snapshot; invoked at proof line 384 of the counted proof. It is retained only as a navigation pointer; every load-bearing step is inside the delivered span.\par\medskip
\ifdefined\hypertarget\hypertarget{sourcepointer-2}{}\fi
\noindent\textit{Source reference.} Original lemma 0A6I: sources/stacks-chow/tags.txt, source lines 12846--12846. Cross-chapter standard prerequisite located in the project tag list staged in the local snapshot; invoked at proof line 393 of the counted proof. It is retained only as a navigation pointer; every load-bearing step is inside the delivered span.\par\medskip
\ifdefined\hypertarget\hypertarget{sourcepointer-3}{}\fi
\noindent\textit{Source reference.} Original lemma 0583: sources/stacks-chow/tags.txt, source lines 6766--6766. Cross-chapter standard prerequisite located in the project tag list staged in the local snapshot; invoked at proof line 404 of the counted proof. It is retained only as a navigation pointer; every load-bearing step is inside the delivered span.\par\medskip
\begin{lemma}
\label{lemma-affine-duality}
Let $A$ be a Noetherian ring and let $X = \Spec(A)$. Let $K, L$ be objects
of $D(A)$. If $K \in D_{\textit{Coh}}(A)$ and $L$ has finite injective
dimension, then
$$
R\SheafHom_{\mathcal{O}_X}(\widetilde{K}, \widetilde{L})
=
\widetilde{R\Hom_A(K, L)}
$$
in $D(\mathcal{O}_X)$.
\end{lemma}

\begin{proof}
We may assume that $L$ is given by a finite complex $I^\bullet$
of injective $A$-modules. By induction on the length of $I^\bullet$
and compatibility of the constructions with distinguished triangles,
we reduce to the case that $L = I[0]$ where $I$ is an injective $A$-module.
In this case, Derived Categories of Schemes, Lemma
\ref{perfect-lemma-quasi-coherence-internal-hom}, tells us that
the $n$th cohomology sheaf of
$R\SheafHom_{\mathcal{O}_X}(\widetilde{K}, \widetilde{L})$
is the sheaf associated to the presheaf
$$
D(f) \longmapsto \Ext^n_{A_f}(K \otimes_A A_f, I \otimes_A A_f)
$$
Since $A$ is Noetherian, the $A_f$-module $I \otimes_A A_f$ is injective
(Dualizing Complexes, Lemma
\ref{dualizing-lemma-localization-injective-modules}). Hence we see that
\begin{align*}
\Ext^n_{A_f}(K \otimes_A A_f, I \otimes_A A_f)
& =
\Hom_{A_f}(H^{-n}(K \otimes_A A_f), I \otimes_A A_f) \\
& =
\Hom_{A_f}(H^{-n}(K) \otimes_A A_f, I \otimes_A A_f) \\
& =
\Hom_A(H^{-n}(K), I) \otimes_A A_f
\end{align*}
The last equality because $H^{-n}(K)$ is a finite $A$-module, see
Algebra, Lemma \ref{algebra-lemma-hom-from-finitely-presented}.
This proves that the canonical map
$$
\widetilde{R\Hom_A(K, L)}
\longrightarrow
R\SheafHom_{\mathcal{O}_X}(\widetilde{K}, \widetilde{L})
$$
is a quasi-isomorphism in this case and the proof is done.
\end{proof}

\end{document}
